% GENERATED FILE. Edit canonical lessons, then run npm run generate:books.
% canonical-source-hash: fnv1a64:74c5c1f90e036cf4
% canonical-lessons: JA-C527-genzai, JA-C527-mirai, JA-C527-saikin, JA-C527-korekara, JA-C527-gendai, JA-R527-review-522-527

\chapter{Now and the Future}
\label{ch:ja-now-and-the-future}
\hlchaptermodality{527}

\begin{chapteropening}
\textbf{By the end of this chapter:} \emph{I can talk about now, recently and what lies ahead.}

Complete the last lesson of chapter 527: I can talk about now, recently and what lies ahead.
\end{chapteropening}

\section[\texorpdfstring{\ja{げんざい} (genzai)}{げんざい (genzai)}]{\ja{げんざい} (genzai) — the present; now}

\label{lesson:JA-C527-genzai}



\begin{quote}
Before the new one: say the Japanese for to turn for the better, then the Japanese for rapidly.
\end{quote}

\subsection*{You'll want to know: \ja{げんざい}}
\textbf{\ja{げんざい}} — \emph{genzai} — \textquotedblleft{}the present; now\textquotedblright{}.

\begin{quote}
\textbf{\ja{げんざいの} \ja{じかんは} \ja{くにに} \ja{よります。}} — \emph{genzai no jikan wa kuni ni yorimasu} — \textquotedblleft{}The time right now depends on the country.\textquotedblright{}
\end{quote}

\begin{cousinweb}
\textbf{\ja{げんざい}} (\emph{genzai}) is Sino-Japanese (kango): two Chinese characters read with their on'yomi, readings borrowed from Chinese. They are \emph{gen}, \textquotedblleft{}present\textquotedblright{} (Middle Chinese \emph{henH}) and \emph{zai}, \textquotedblleft{}exist\textquotedblright{} (\emph{dzojX}) (KANJIDIC2; Middle Chinese in Baxter's transcription).
\end{cousinweb}

\begin{grammarlens}[title={What you've built}]
One more word for saying what happens if.
\end{grammarlens}

\subsection*{Your turn}
\begin{itemize}
\raggedright
  \item \emph{Say it:} \emph{genzai}
  \item \emph{Say it:} \emph{genzai}, once more
  \item \emph{Say it:} \emph{genzai no jikan wa kuni ni yorimasu} — The time right now depends on the country.
  \item \emph{Recall:} say the Japanese for to turn for the better, then the Japanese for rapidly, then say \emph{genzai} again
\end{itemize}

\begin{tcolorbox}[breakable,colback=teal!4,colframe=teal!35!black,title={Before you move on}]
What does \ja{げんざい} mean? (\textquotedblleft{}The present; now\textquotedblright{}.) Say it once more.
\end{tcolorbox}
\section[\texorpdfstring{\ja{みらい} (mirai)}{みらい (mirai)}]{\ja{みらい} (mirai) — the distant future}

\label{lesson:JA-C527-mirai}



\begin{quote}
Before the new one: say the Japanese for rapidly, then the Japanese for the present.
\end{quote}

\subsection*{You'll want to know: \ja{みらい}}
\textbf{\ja{みらい}} — \emph{mirai} — \textquotedblleft{}the distant future\textquotedblright{}.

\begin{quote}
\textbf{\ja{みらいは} \ja{いまの} \ja{まいにちに} \ja{よります。}} — \emph{mirai wa ima no mainichi ni yorimasu} — \textquotedblleft{}The future depends on how we live each day now.\textquotedblright{}
\end{quote}

\begin{cousinweb}
\textbf{\ja{みらい}} (\emph{mirai}) is Sino-Japanese (kango): two Chinese characters read with their on'yomi, readings borrowed from Chinese. They are \emph{mi}, \textquotedblleft{}not yet\textquotedblright{} and \emph{rai}, \textquotedblleft{}come\textquotedblright{} (Middle Chinese \emph{loj}) (KANJIDIC2; Middle Chinese in Baxter's transcription).
\end{cousinweb}

\begin{grammarlens}[title={What you've built}]
One more word for saying what happens if.
\end{grammarlens}

\subsection*{Your turn}
\begin{itemize}
\raggedright
  \item \emph{Say it:} \emph{mirai}
  \item \emph{Say it:} \emph{mirai}, once more
  \item \emph{Say it:} \emph{mirai wa ima no mainichi ni yorimasu} — The future depends on how we live each day now.
  \item \emph{Recall:} say the Japanese for rapidly, then the Japanese for the present, then say \emph{mirai} again
\end{itemize}

\begin{tcolorbox}[breakable,colback=teal!4,colframe=teal!35!black,title={Before you move on}]
What does \ja{みらい} mean? (\textquotedblleft{}The distant future\textquotedblright{}.) Say it once more.
\end{tcolorbox}
\section[\texorpdfstring{\ja{さいきん} (saikin)}{さいきん (saikin)}]{\ja{さいきん} (saikin) — recently; these days}

\label{lesson:JA-C527-saikin}



\begin{quote}
Before the new one: say the Japanese for the present, then the Japanese for the distant future.
\end{quote}

\subsection*{You'll want to know: \ja{さいきん}}
\textbf{\ja{さいきん}} — \emph{saikin} — \textquotedblleft{}recently; these days\textquotedblright{}.

\begin{quote}
\textbf{\ja{さいきんの} \ja{てんきは} \ja{ばしょに} \ja{よります。}} — \emph{saikin no tenki wa basho ni yorimasu} — \textquotedblleft{}The weather lately depends on where you are.\textquotedblright{}
\end{quote}

\begin{cousinweb}
\textbf{\ja{さいきん}} (\emph{saikin}) is Sino-Japanese (kango): two Chinese characters read with their on'yomi, readings borrowed from Chinese. They are \emph{sai}, \textquotedblleft{}utmost\textquotedblright{} (Middle Chinese \emph{tswajH}) and \emph{kin}, \textquotedblleft{}near\textquotedblright{} (\emph{gj+nX}) (KANJIDIC2; Middle Chinese in Baxter's transcription).
\end{cousinweb}

\begin{grammarlens}[title={What you've built}]
One more word for saying what happens if.
\end{grammarlens}

\subsection*{Your turn}
\begin{itemize}
\raggedright
  \item \emph{Say it:} \emph{saikin}
  \item \emph{Say it:} \emph{saikin}, once more
  \item \emph{Say it:} \emph{saikin no tenki wa basho ni yorimasu} — The weather lately depends on where you are.
  \item \emph{Recall:} say the Japanese for the present, then the Japanese for the distant future, then say \emph{saikin} again
\end{itemize}

\begin{tcolorbox}[breakable,colback=teal!4,colframe=teal!35!black,title={Before you move on}]
What does \ja{さいきん} mean? (\textquotedblleft{}Recently; these days\textquotedblright{}.) Say it once more.
\end{tcolorbox}
\section[\texorpdfstring{\ja{これから} (korekara)}{これから (korekara)}]{\ja{これから} (korekara) — from now on; after this}

\label{lesson:JA-C527-korekara}



\begin{quote}
Before the new one: say the Japanese for the distant future, then the Japanese for recently.
\end{quote}

\subsection*{You'll want to know: \ja{これから}}
\textbf{\ja{これから}} — \emph{korekara} — \textquotedblleft{}from now on; after this\textquotedblright{}.

\begin{quote}
\textbf{\ja{これから} \ja{さむく} \ja{なったら、かぜに} \ja{ちゅういして} \ja{ください。}} — \emph{korekara samuku nattara, kaze ni chūi shite kudasai} — \textquotedblleft{}If it gets colder from now on, take care not to catch a cold.\textquotedblright{}
\end{quote}

\begin{cousinweb}
\textbf{\ja{これから}} (\emph{korekara}) is a native Japanese word (wago), inherited rather than borrowed from Chinese, and usually written in kana; its older origin is not settled (JMdict).
\end{cousinweb}

\begin{grammarlens}[title={What you've built}]
One more word for saying what happens if.
\end{grammarlens}

\subsection*{Your turn}
\begin{itemize}
\raggedright
  \item \emph{Say it:} \emph{korekara}
  \item \emph{Say it:} \emph{korekara}, once more
  \item \emph{Say it:} \emph{korekara samuku nattara, kaze ni chūi shite kudasai} — If it gets colder from now on, take care not to catch a cold.
  \item \emph{Recall:} say the Japanese for the distant future, then the Japanese for recently, then say \emph{korekara} again
\end{itemize}

\begin{tcolorbox}[breakable,colback=teal!4,colframe=teal!35!black,title={Before you move on}]
What does \ja{これから} mean? (\textquotedblleft{}From now on; after this\textquotedblright{}.) Say it once more.
\end{tcolorbox}
\section[\texorpdfstring{\ja{げんだい} (gendai)}{げんだい (gendai)}]{\ja{げんだい} (gendai) — the present day; modern times}

\label{lesson:JA-C527-gendai}



\begin{quote}
Before the new one: say the Japanese for recently, then the Japanese for from now on.
\end{quote}

\subsection*{You'll want to know: \ja{げんだい}}
\textbf{\ja{げんだい}} — \emph{gendai} — \textquotedblleft{}the present day; modern times\textquotedblright{}.

\begin{quote}
\textbf{\ja{もし} \ja{そふの} \ja{そふが} \ja{げんだいに} \ja{きたら、おどろくでしょう。}} — \emph{moshi sofu no sofu ga gendai ni kitara, odoroku deshō} — \textquotedblleft{}If my grandfather's grandfather came to the present day, he would be amazed.\textquotedblright{}
\end{quote}

\begin{cousinweb}
\textbf{\ja{げんだい}} (\emph{gendai}) is Sino-Japanese (kango): two Chinese characters read with their on'yomi, readings borrowed from Chinese. They are \emph{gen}, \textquotedblleft{}present\textquotedblright{} (Middle Chinese \emph{henH}) and \emph{dai}, \textquotedblleft{}age\textquotedblright{} (\emph{dojH}) (KANJIDIC2; Middle Chinese in Baxter's transcription).
\end{cousinweb}

\begin{grammarlens}[title={What you've built}]
One more word for saying what happens if.
\end{grammarlens}

\subsection*{Your turn}
\begin{itemize}
\raggedright
  \item \emph{Say it:} \emph{gendai}
  \item \emph{Say it:} \emph{gendai}, once more
  \item \emph{Say it:} \emph{moshi sofu no sofu ga gendai ni kitara, odoroku deshō} — If my grandfather's grandfather came to the present day, he would be amazed.
  \item \emph{Recall:} say the Japanese for recently, then the Japanese for from now on, then say \emph{gendai} again
\end{itemize}

\begin{tcolorbox}[breakable,colback=teal!4,colframe=teal!35!black,title={Before you move on}]
What does \ja{げんだい} mean? (\textquotedblleft{}The present day; modern times\textquotedblright{}.) Say it once more.
\end{tcolorbox}
\section[(dialogue)]{Looking back: from Something, Someone, Somewhere to Now and the Future}

\label{lesson:JA-R527-review-522-527}



\begin{quote}
Say the Japanese for something and someone.
\end{quote}

\subsection*{Your turn}
Say these aloud:

\begin{itemize}
\raggedright
  \item something, someone, somewhere, at any time, anything
  \item anyone, any one (of them), no doubt, as it happens, somehow (or other)
  \item a supposition, to picture (in one's mind), to imagine, reasoning, to be realized
  \item the real thing, a fake, (not) in the least, to compare (to), sooner or later
  \item to progress, to develop, to improve (something), to turn for the better, rapidly
  \item the present, the distant future, recently, from now on, the present day
\end{itemize}

\begin{tcolorbox}[breakable,colback=teal!4,colframe=teal!35!black,title={Before you move on}]
Something? (\textbf{\ja{なにか}}, \emph{nanika}.) Someone? (\textbf{\ja{だれか}}, \emph{dareka}.) Somewhere? (\textbf{\ja{どこか}}, \emph{dokoka}.) At any time? (\textbf{\ja{いつでも}}, \emph{itsudemo}.) Anything? (\textbf{\ja{なんでも}}, \emph{nandemo}.) Anyone? (\textbf{\ja{だれでも}}, \emph{daredemo}.) Any one (of them)? (\textbf{\ja{どれでも}}, \emph{doredemo}.) No doubt? (\textbf{\ja{さぞ}}, \emph{sazo}.) As it happens? (\textbf{\ja{たまたま}}, \emph{tamatama}.) Somehow (or other)? (\textbf{\ja{なんとか}}, \emph{nantoka}.) A supposition? (\textbf{\ja{かてい}}, \emph{katei}.) To picture (in one's mind)? (\textbf{\ja{おもいえがく}}, \emph{omoiegaku}.) To imagine? (\textbf{\ja{そうぞうする}}, \emph{sōzō suru}.) Reasoning? (\textbf{\ja{すいり}}, \emph{suiri}.) To be realized? (\textbf{\ja{じつげんする}}, \emph{jitsugen suru}.) The real thing? (\textbf{\ja{ほんもの}}, \emph{honmono}.) A fake? (\textbf{\ja{にせもの}}, \emph{nisemono}.) (not) in the least? (\textbf{\ja{まるで}}, \emph{marude}.) To compare (to)? (\textbf{\ja{たとえる}}, \emph{tatoeru}.) Sooner or later? (\textbf{\ja{いずれ}}, \emph{izure}.) To progress? (\textbf{\ja{しんぽする}}, \emph{shinpo suru}.) To develop? (\textbf{\ja{はってんする}}, \emph{hatten suru}.) To improve (something)? (\textbf{\ja{かいぜんする}}, \emph{kaizen suru}.) To turn for the better? (\textbf{\ja{こうてんする}}, \emph{kōten suru}.) Rapidly? (\textbf{\ja{どんどん}}, \emph{dondon}.) The present? (\textbf{\ja{げんざい}}, \emph{genzai}.) The distant future? (\textbf{\ja{みらい}}, \emph{mirai}.) Recently? (\textbf{\ja{さいきん}}, \emph{saikin}.) From now on? (\textbf{\ja{これから}}, \emph{korekara}.) The present day? (\textbf{\ja{げんだい}}, \emph{gendai}.)
\end{tcolorbox}
